\documentclass[UTF8]{article}
\usepackage[utf8]{inputenc}
\usepackage{indentfirst}
\usepackage{graphicx}
\usepackage{hyperref}
\usepackage{ctex}
\usepackage{geometry}
\geometry{a4paper,left=2.5cm,right=2.5cm,top=2.5cm,bottom=2.5cm}

\hypersetup{
    colorlinks=true,
    linkcolor=blue,
    urlcolor=blue
}

\title{个人主页搭建报告}
\author{田天威 PB23010364}
\date{}

\begin{document}

\maketitle

\section{个人主页介绍}

\subsection{主页访问链接}
个人主页已部署至 GitHub Pages，访问地址为：\url{https://Yomeko1.github.io/Personal-Page/}

\subsection{主页的主要内容}
个人主页包含以下内容：
\begin{itemize}
    \item \textbf{标签}：包含姓名、学校与专业方向（USTC | C/C++ | Python | 信息与计算科学）
    \item \textbf{个人简介}：介绍个人背景、兴趣爱好及学习方向
    \item \textbf{研究方向与兴趣}：涵盖信息与计算科学、计算机图形学、人工智能等领域
    \item \textbf{学习笔记}：提供微分几何学习笔记的 PDF 
    \item \textbf{项目展示}：预留项目展示区域（目前为 "coming soon"）
    \item \textbf{博客入口}：提供进入技术博客的跳转链接
\end{itemize}

\subsection{页面结构与功能说明}
主页采用单页结构，通过垂直滚动浏览全部内容。页面顶部设有导航栏，包含"首页"和"博客"两个入口，便于用户在不同页面间切换。页面采用响应式布局设计，能够适配不同屏幕尺寸的设备。

页面整体使用紫色渐变色作为主题色调，视觉风格统一、简洁美观。

\section{博客内容介绍}

\subsection{博客主题}
本次技术博客的主题为\textbf{计算机图形学}，具体探讨了游戏中常见的光线追踪（Ray Tracing）技术。

博文标题为《从光栅到光追：游戏显卡的算力消耗解析》。

\subsection{主要内容概述}
博客从游戏设置中的“光线追踪”选项出发，系统性地介绍了以下内容：

\begin{itemize}
    \item \textbf{光栅化渲染（Rasterization）}：传统渲染方式的原理与 GPU 并行架构
    \item \textbf{光线追踪（Ray Tracing）}：光线方程、与球体/三角形求交的数学原理
    \item \textbf{路径追踪与俄罗斯轮盘赌}：蒙特卡罗方法、大数定律收敛、无偏终止策略
    \item \textbf{性能对比分析}：解释为何开启光线追踪后电脑会出现卡顿现象
    \item \textbf{现代游戏混合方案}：光栅化 + 有限光追 + AI 降噪的工程实践
\end{itemize}

博客中包含了详细的数学公式推导（如光线参数方程 \(P(t) = O + t \cdot d\)、球体求交的二次方程、反射与折射公式等），并配有部分示意图。

\subsection{博客访问方式}
博客可通过个人主页顶部的"博客"导航栏进入，访问链接为：\url{https://Yomeko1.github.io/Personal-Page/blog.html}

光线追踪专题文章直接访问：\url{https://Yomeko1.github.io/Personal-Page/blog-ray-tracing.html}

\section{实现过程}

\subsection{使用的主要工具或技术}

\begin{itemize}
    \item \textbf{HTML/CSS}：使用纯 HTML 和 CSS 构建静态网页，零依赖
    \item \textbf{Git 版本管理}：使用 Git 管理项目代码，实现版本控制
    \item \textbf{GitHub Pages}：利用 GitHub Pages 进行网页托管与部署
    \item \textbf{MathJax}：引入 MathJax 实现 LaTeX 数学公式的网页渲染
    \item \textbf{AI 辅助工具}：使用 AI Agent（OpenCode）辅助搭建个人主页与撰写博客
\end{itemize}

\subsection{网页搭建和部署过程}

网页搭建与部署的完整流程如下：

\begin{enumerate}
    \item \textbf{项目初始化}：创建本地 Git 仓库，初始化项目结构
    \item \textbf{页面开发}：编写 index.html（主页）、blog.html（博客列表）、blog-ray-tracing.html（博客正文）以及 styles.css（样式文件）
    \item \textbf{内容填充}：填充个人资料，包括头像、姓名、学校、兴趣方向等
    \item \textbf{笔记与项目集成}：将微分几何学习笔记（main.pdf）链接到主页笔记展示区，项目展示区暂定 Coming Soon
    \item \textbf{提交与推送}：使用 Git 提交代码，通过 \texttt{git push} 推送到 GitHub 远程仓库
    \item \textbf{GitHub Pages 配置}：在 GitHub 仓库 Settings 中启用 Pages 功能，选择 master 分支作为部署源
\end{enumerate}

\subsection{遇到的问题及解决方法}

\begin{itemize}
    \item \textbf{网络连接问题}：在推送代码至 GitHub 时遇到网络连接失败，经过多次重试后成功推送
    \item \textbf{图片路径问题}：初始头像路径使用本地文件名，后统一将图片资源移至 Image 文件夹，修改 HTML 中的相对路径
    \item \textbf{图片显示尺寸问题}：部分图片（扫描线算法、BVH 树）显示尺寸过大，通过 CSS 设置 \texttt{max-width} 属性限制最大宽度
    \item \textbf{跨客户端显示不一致问题}：网页在手机端存在图片、公式超出边界的情况，通过添加响应式 CSS 规则（如 \texttt{overflow-x: auto}）使页面在桌面端和移动端均能正常显示
\end{itemize}

\section{AI Agent 使用说明}

\subsection{使用的 AI 工具}
本次大作业使用 OpenCode 作为 AI Agent 辅助工具，进行网页搭建与博客撰写。

\subsection{AI 辅助完成的内容}

\begin{itemize}
    \item \textbf{网页框架搭建}：AI 辅助生成了主页、博客列表页、博客内容页的基本 HTML 结构与 CSS 样式，并完成了基本信息的占位。
    \item \textbf{博客内容撰写}：AI 辅助撰写了光线追踪技术博客，包括数学公式（如光线参数方程、球体求交公式、路径追踪积分等）和算法伪代码（递归光线追踪、俄罗斯轮盘赌）
    \item \textbf{技术细节补充}：AI 补充了路径追踪的蒙特卡罗采样方法、俄罗斯轮盘赌的无偏终止策略等技术内容
    \item \textbf{图片与排版优化}：AI 协助调整图片显示尺寸、添加图片版权声明、配置 MathJax 公式渲染
\end{itemize}

\subsection{自己对 AI 生成内容做的修改}

\begin{itemize}
    \item \textbf{内容个性化}：将 AI 生成的初始内容中的占位信息替换为个人真实信息，简介文字也根据个人实际情况进行了修改
    \item \textbf{主题选择}：博客主题从最初的数学主题调整为计算机图形学中的光线追踪内容，更符合个人兴趣与研究方向
    \item \textbf{修改内容}：增加了个人的学习笔记，并对博客内容进行了一部分修改
    \item \textbf{细节调整}：对 AI 生成的公式表述、代码注释进行了检查与微调，确保数学公式的正确性与文字表达的流畅性
\end{itemize}

\subsection{对 AI 辅助编程的思考}

通过本次大作业，我深刻体会到 AI Agent 在辅助编程方面的强大能力：

\begin{itemize}
    \item \textbf{效率提升}：AI 能够快速生成网页框架与样式代码，大幅减少了重复性编码工作的时间
    \item \textbf{知识拓展}：在撰写光线追踪博客的过程中，AI 能够提供系统的知识梳理与准确的数学公式，降低了学习门槛
    \item \textbf{协作流程}：使用 AI 辅助编程时，我倾向于先让 AI 生成初稿，再根据个人需求进行个性化修改，这种人机协作模式既保证了效率，又确保了内容的真实性
    \item \textbf{质量把控}：尽管 AI 生成的内容在结构和语法上通常没有问题，但仍需人工检查内容的准确性与适用性，特别是在专业性较强的技术博客中
\end{itemize}

总体而言，AI 辅助编程是一个强有力的工具，能够帮助我们更高效地完成项目搭建与内容创作。关键在于如何合理地利用 AI 的能力，同时保持对最终成果的质量把控。

\section{结果展示}

\subsection{个人主页截图}
个人主页采用紫色渐变主题，顶部展示头像、姓名与标签信息，下方依次为"关于我"、"研究方向与兴趣"、"学习笔记"、"项目展示"等模块，底部设有导航栏与博客入口。

主页访问地址：\url{https://Yomeko1.github.io/Personal-Page/}

\subsection{博客页面截图}
博客列表页展示技术博客入口，点击可进入光线追踪专题文章。博客正文包含完整的数学公式、算法伪代码与示意图，支持点击目录跳转功能。

博客访问地址：\url{https://Yomeko1.github.io/Personal-Page/blog.html}

文章地址：\url{https://Yomeko1.github.io/Personal-Page/blog-ray-tracing.html}

\subsection{其他说明}
整个项目已通过 Git 进行版本管理，代码存放于 GitHub 仓库：\url{https://github.com/Yomeko1/Personal-Page}

所有页面均支持响应式布局，能够在桌面端与移动端设备上正常显示。

\end{document}